% Introduction. Labels defined here use the prefix "intro:".
\section{Introduction}\label{sec:introduction}

Branch-and-bound produces a feasible solution and a proof of its quality.
The proof excludes every region that might contain a substantially better
solution, using relaxation bounds, disjunctions, and domain reductions. Its
size can vary greatly between instances of the same dimension. A large
search tree can reflect a weak relaxation, poorly chosen splits, expensive
propagation, or difficulty finding a good incumbent. These causes require
different explanations.

This paper studies the size of the proof for a fixed instance, relaxation,
and class of permitted operations. We call this size the
\emph{certification complexity}. The central question is how the
relaxation's error interacts with the instance's near-optimal points.
We develop geometric characterizations for continuous box relaxations,
combinatorial lower bounds for convex integer relaxations, and probabilistic
laws in two Gaussian models. We then examine how propagation, branching,
and decomposition change the certificates that can be constructed.

\subsection{Certificate size and search work}\label{intro:certificates}

Let $f_*$ be the global minimum and let $\epsilon>0$ be an additive
optimality tolerance. A lower-bound certificate covers the points that must
be considered by regions on which a specified test proves that the objective
is at least $f_*-\epsilon$. A finished branch-and-bound tree supplies such a
certificate when its incumbent is optimal. With any feasible incumbent of
value $U\geq f_*$, pruning at $L(B)\geq U-\epsilon$ still implies
$L(B)\geq f_*-\epsilon$, so the same lower bounds apply. Constructive upper
bounds usually fix an optimal incumbent; the relevant exceptions state an
explicit search convention.

For spatial certificates, three minima are useful:
$N_{\mathrm{cover}}$ counts a smallest cover by valid boxes,
$N_{\mathrm{rect}}$ a smallest rectangular partition, and
$N_{\mathrm{tree}}$ the leaves of a smallest tree formed by coordinate
cuts. They satisfy
\[
 N_{\mathrm{cover}}\leq N_{\mathrm{rect}}\leq N_{\mathrm{tree}}.
\]
These finite minima need not be equal. A lower bound for covers and an upper
bound realized by a tree can nevertheless identify their common asymptotic
order. A complete binary tree with $L$ leaves has $2L-1$ nodes.

Domain reduction requires additional care. A closed discarded slab can
contain a retained feasible face, whose points the reduction did not
exclude. We therefore pair each closed test box with the subset of points
it \emph{owns}; retained boundaries continue into the surviving region.
These owned families need not be ordinary covers by boxes valid on their
whole closures. The geometric lower bounds apply to them directly through
per-test estimates. \Cref{sec:certificates} gives the ownership construction
and charges removed pieces, probe tests used in stopping bounds, and
relaxation evaluations (\cref{cert:lem-events}). Lower bounds after reduction require the stated
pointwise gap assumptions. A weaker bound-gap assumption suffices for raw
certificates without reduction.

We distinguish four quantities:
\begin{enumerate}[label=(\roman*)]
\item the size of a smallest certificate in the chosen class;
\item the number of nodes processed while finding one;
\item the work needed to check its regions, including relaxation evaluations,
  tightening rounds, and propagation rounds;
\item the work needed to find an incumbent of the required quality.
\end{enumerate}
The branching results compare item~(ii) with item~(i). Propagation and
probing can shift work from nodes to item~(iii). A geometric lower bound on
owned events becomes a node lower bound only when the relevant rounds per
node are bounded and the model accounts for probe tests. Every result fixes
its relaxation and operation class; a stronger relaxation or an additional
operation must satisfy the hypotheses anew.

\subsection{Main results}\label{intro:results}

The continuous results concern $\epsilon\downarrow0$ for a fixed instance.
Write $D\subset\R^n$ for the compact root box, $F\subseteq D$ for its
feasible set, $m=f-f_*$, and
$E(\eta)=\{x\in F:m(x)\leq\eta\}$. Constants may depend on the instance,
relaxation, and dimension. The random-model results instead let the problem
size grow and state their probability and scaling assumptions separately.

\paragraph{Near-optimal geometry.}
For a box $B=\prod_i[l_i,u_i]$, put
\[
 q_B(x)=\sum_i(x_i-l_i)(u_i-x_i).
\]
Suppose the bound satisfies $L(B)\leq f(x)-\alpha q_B(x)$ at every
feasible $x\in B$, for a fixed $\alpha>0$. A valid box then obeys
$\alpha q_B(x)\leq m(x)+\epsilon$. Each summand of $q_B$ is at least
the squared distance to the nearer endpoint, so every point of
$E(\eta)$ certified by that box lies within sup-norm distance
$\sqrt{(\epsilon+\eta)/\alpha}$ of a box vertex. This gives the
covering lower bound that drives the continuous theory.

Let $N_\infty(S,r)$ count the least number of axis-aligned cubes of side
$r$ covering $S$, and define
\begin{equation}\label{intro:eq-profile}
 \Phi_\alpha(\epsilon)=\sup_{\eta\geq0}
 N_\infty\!\left(E(\eta),2\sqrt{(\epsilon+\eta)/\alpha}\right).
\end{equation}
Under second-order upper error in both objective and constraint violation,
a local violation-to-distance error bound near the full feasible set, and a
Lipschitz objective, dyadic subdivision yields
\begin{equation}\label{intro:eq-comparison}
 c\Phi_\alpha(\epsilon)
 \leq N_{\mathrm{cover}}\leq N_{\mathrm{rect}}\leq N_{\mathrm{tree}}
 \leq T_{\mathrm{bis}}
 \leq C_0+C\log(1/\epsilon)\Phi_\alpha(\epsilon)
\end{equation}
for sufficiently small $\epsilon$ (\cref{geom:thm-profile}). The supremum over
levels is essential: one near-optimal level can miss many separated wells.
With quadratic growth away from the optimal set and the specified density
conditions, positive-dimensional optimal manifolds give the familiar
$\epsilon^{-p/2}$ rate for dimension $p$. An isolated degenerate minimizer
can also give a polynomial rate, so the dimension of the optimal set alone
does not determine complexity.

Integral lower bounds sharpen this geometry on strata with controlled
coordinate-projection multiplicity. Relaxed constraints can supply further
witnesses: slightly infeasible points with objective below $f_*-\epsilon$
must be excluded too. Under an inner-tube hypothesis, a global KKT minimizer
with a positive-dimensional active stratum and a nonzero multiplier for a
relaxed constraint forces $\Omega(\log(1/\epsilon))$ all-owner events. This
result excludes reductions using original feasibility. It allows nonlinear
exactly kept constraints, with a $C^2$ objective and compatible original
exact and relaxed constraint data, and LICQ and multipliers for that same
description. Changing only the description
of the feasible set does not preserve these hypotheses (\cref{geom:thm-tube-kkt}).

\paragraph{All root-box faces can matter.}
For unconstrained $f$ with a Lipschitz gradient, a lower bound gap
$\alpha q_B$, and pointwise upper error $\alpha' q_B$, the three spatial
minima and dyadic node count are within constant factors of
\begin{equation}\label{intro:eq-faces}
 1+\sum_{G:\,\dim G\geq1}
 \int_G(m(x)+\epsilon)^{-\dim G/2}\,d\mathcal H^{\dim G}(x),
\end{equation}
where $G$ ranges over the faces of $D$, including $D$ itself
(\cref{face:thm:characterization,face:rem:gap-scope}). This characterization
needs no quadratic-growth
assumption. For analytic objectives, classical analytic sublevel asymptotics
convert each face integral into a polynomial or logarithmic rate. On a
nonzero face restriction with zeros, small sublevel volumes have order
$t^{\lambda_G}\log^{\theta_G-1}(1/t)$; $\lambda_G$ is the real log
canonical threshold and $\theta_G$ its multiplicity. That face contributes
$\epsilon^{\lambda_G-\dim G/2}\log^{\theta_G-1}(1/\epsilon)$ when
$\lambda_G<\dim G/2$, and $\log^{\theta_G}(1/\epsilon)$ at equality;
positive-minimum and identically-zero faces are treated separately.
When all minimizers are interior, the full-dimensional integral suffices,
and the rate is logarithmic exactly when all minimizers are nondegenerate.
At the boundary a proper face can give a larger exponent than the full box
(\cref{face:cor:analytic-rate,face:ex:boundary}).

McCormick relaxations have a different exactness geometry. For separately
relaxed bilinear terms, the zero-gap faces correspond to vertex covers of
the interaction graph. \Cref{sec:geometry} derives matching, transversality,
and fractional-cover bounds from this structure. A two-dimensional convex
example with an aligned optimal segment has a two-leaf exact certificate
(\cref{geom:prop-aligned-two}), while an aligned three-dimensional example
still requires
$\Theta(\epsilon^{-1/2})$ boxes (\cref{geom:ex-aligned-three}). Alignment
alone therefore does not supply a
general characterization. These results explain why convergence order alone
can miss the placement of exact regions that determines certificate size.

\paragraph{Changing the certificate representation.}
\Cref{sec:decomposition} retains a fixed factorization and local relaxation,
but allows local bag certificates connected by affine separator minorants.
For a tree decomposition with $M$ bags and width $w$, global quadratic
growth, $C^{1,1}$ bag objectives, vertex-vanishing factor error, and fixed width,
incidence, and conditioning parameters give a certificate of size
$O(M\log(M/\epsilon))$, with explicit dependence on those parameters
(\cref{decomp:small-certificate}). This is an existence result using the
minimizer and separator slopes. A nonconvex quartic path family has a unique
nondegenerate minimizer and an exponential spatial-cover lower bound in
dimension, yet admits an $O(n\log(n/\epsilon))$ decomposition certificate
under the same separate unary and bilinear oracle
(\cref{decomp:path-separation}). The comparison concerns
proof size. Construction, local checking queries, and solution time carry
separate costs.

\paragraph{Propagation and branching.}
\Cref{sec:propagation} models objective-only cutoff propagation by the
greatest hull-consistent fixed point of a specified finite expression graph.
Local exactness gives an ideal certificate count independent of tolerance;
one-sided graphs supply a sufficient condition. Conversely, specified
coordinatewise derivative cancellation with a uniform no-dominance margin
transfers geometric lower bounds through charged propagation phases. These
are sufficient conditions in different classes. The round count remains
schedule dependent, and a bounded ideal node count can coexist with a
diverging number of rounds
(\cref{prop:local-exactness,prop:cnd-transfer}).

For the exact oracle $L(B)=\min_B(f-\alpha q_B)$ on an interval,
splitting at any relaxation minimizer is 4-competitive against the smallest
finished tree, with an optimal incumbent and continuous $f$
(\cref{branching:four-competitive}). Convexity is unnecessary for this
counting theorem.
Smooth nonanalytic examples give information lower bounds for node-local
rules, and fixed interior safeguards impose an unbounded worst-case loss on
kink families. A recentring safeguard attains the shortest safe chain that
hits the kink; its finite-tolerance bound is stated separately. In several
dimensions, splitting every available coordinate can lose a logarithmic
factor even on a separable quadratic. Repeated coordinate ambiguity forces
an exponential dependence on dimension for node-local rules with
coordinatewise minimizer selection, nonanalytic germ information, and no
objective-dependent information from siblings
(\cref{branching:dimension-obstruction}). The rule selecting the largest
coordinate contribution to $q_B$ at the relaxation minimizer has a constant
ratio for one arbitrary separable coordinate and a fixed number of sharp
coordinates, with an exact coordinatewise oracle satisfying the sharpness
assumptions (\cref{branching:sharp-induction}). Its general
fixed-dimensional competitiveness remains open.

\paragraph{Integer certificates.}
For a fixed convex projected relaxation $\phi$, a required integer set
$P$, and threshold $\tau$, call a class admissible if $\phi\geq\tau$
on its convex hull. The minimum number of such classes covering $P$ is the
\emph{class number} $\kappa_\tau(P)$. It bounds the leaves of every
$P$-covering convex-piece certificate. An unrestricted semantic binary
certificate attains $\kappa_\tau(P)$ leaves when its root may be any convex
set containing $P$ and its convex children cover the required integer
points. This existence statement asserts no efficient construction or
checking algorithm (\cref{lattice:class-number}). Convex cuts preserving
every required point in $P$ preserve the lower bound. Certified incumbent removals require an augmented
leaf-and-removal count.

The class number can be much smaller than the best variable-branching
tree. For a Haar random unimodular lattice with a uniform torus target, vanishing tolerance relative
to the Gaussian-heuristic squared radius gives class-number bounds between
$2^{(\frac12\log_2(3/2)-o(1))n}$ and $2^{(1/2+o(1))n}$ with high probability.
They are invariant under basis changes (\cref{lattice:cvp}). A separate
block-curvature family
has exponentially many classes for a perspective relaxation and an exact
root bound from small block hulls.

\paragraph{Gaussian statistical models.}
Condition C1 requires every single-variable fixing that disagrees with a
specified optimal binary point to have relaxation bound at least $\OPT$.
When the relaxation is exact at fully fixed assignments, it gives a path
certificate with at most $2p+1$ nodes for $p$ binary variables if the optimal
incumbent is available. Strict bounds above $\OPT$ give the same search
bound under exact best-bound selection when a fully fixed optimal node
recovers its feasible extension (\cref{lattice:path}). Failure
of C1 alone gives no lower bound on tree size.

For sparse ridge regression, \cref{sec:regression} identifies separate
planted-support thresholds for root and single-variable certification.
They are governed by $2\log p$ and $2\log(p\lambda/n)$ for an explicit
ridge-dependent statistic, under the stated Gaussian high-dimensional
scaling and fixed margins. With $k=p^{\gamma+o(1)}$ and
$n=a k\log p$, a suitable ridge yields a unique planted optimum, an inexact
root, and a linear path certificate when $2-\gamma<a<2$ (\cref{regression:window}).
Selected local second-moment lifts preserve these first-order thresholds
under explicit sandwich and growth assumptions. In a separate pure-noise
or low-total-signal regime, conflict packings force
$\binom pk^{c+o(1)}$ leaves. The pairwise-lift hard bound requires that
zero-fixed helper columns remain in the lift.

For binary least squares, let $A=\sqrt{\rho/N}H$ with $H$ an $M\times N$
standard Gaussian matrix, independent standard Gaussian noise, and
$\beta_N=M/N\to\beta\geq1$. For $\rho\to\infty$, $\rho=o(N)$, and the
tolerances $0\leq\epsilon_N\leq\rho_N$, the logarithm of the minimum convex-piece
certificate size and of the minimum variable-tree leaf count is
$\Theta((N/\rho)\log\rho)$; sufficiently large fixed $\rho$ gives
$\Theta(N)$ instead (\cref{bls:certificate-law}). A finite Gaussian
nonnegative least-squares coefficient bound proves the simultaneous C1
threshold $\theta_c=1/[4(2\beta-1)]$ when $\rho=\theta N$ with fixed
$\theta>0$ separated from $\theta_c$ (\cref{bls:c1-threshold}). It also yields an exact
coefficient mixture and a random root threshold asymptotic to
$\log N/(2\beta_N-1)$ for equality of box and orthant values (\cref{bls:root-threshold}). This root
event concerns upper-box inactivity. Planted recovery, rounding correctness,
and an integral root optimum are distinct events.

\begin{table}[t]
\centering
\small
\setlength{\tabcolsep}{4pt}
\renewcommand{\arraystretch}{1.12}
\begin{tabular}{@{}p{0.22\textwidth}p{0.44\textwidth}p{0.28\textwidth}@{}}
\toprule
Certificate or search model & Governing structure & Scope of the comparison \\
\midrule
Spatial box certificates & Near-optimal covering profile; all root-face
integrals & Specified lower gap and upper-error assumptions \\
McCormick certificates & Interaction graph and placement of near-optimal
strata & Separate lower bounds and exact examples \\
Decomposition certificates & Bag width, incidence, and quadratic growth &
Fixed local oracle; size and query costs distinguished \\
Cutoff propagation & Expression graph and hull-consistent fixed point &
Objective-only phases; rounds counted separately \\
Branching search & Relaxation minimizers and available node information &
Exact-gap oracle; incumbent and selection hypotheses \\
Convex integer certificates & Admissible classes of required integer points &
Fixed projected relaxation and retained point set \\
Gaussian models & Correlation extremes, convex conflicts, and NNLS
coefficients & Distinct asymptotic regimes and certificate events \\
\bottomrule
\end{tabular}
\caption{The certificate models compared in this paper. Each chapter states
its hypotheses and cost measure; a bound in one row does not automatically
transfer to another.}
\label{intro:tab-models}
\end{table}

\subsection{Relation to prior work}\label{intro:prior}

The cluster problem in global optimization concerns the accumulation of
boxes near minimizers. Earlier analyses relate this accumulation to the
convergence order and constants of bounding schemes
\citep{duKearfott1994Cluster,wechsungSchaberBarton2014Cluster}.
Constrained analyses also identify nearly feasible points with low objective
values as a source of clustering
\citep{kannanBarton2017ConstrainedCluster}.
Our profile comparison and root-face characterization address the size of
adaptive certificates, with lower bounds applying to every valid cover and
constructive upper bounds realized by subdivision. The nonlinear
constraint-gap argument retains the original exact/relaxed partition and
makes the operation restrictions explicit.

Instance-dependent certified Lipschitz optimization also uses near-optimal
covering numbers
\citep{bachocCesariGerchinovitz2021LipschitzCertificates}.
That work measures function evaluations under a Lipschitz oracle. Here the
local relaxation specifies the regions that can be certified and the
relevant resolution. Classical analytic singularity theory supplies the
sublevel-volume asymptotics used in \cref{sec:face-exact}
\citep{lin2017LearningCoefficient}; our result connects these asymptotics on
all box faces to spatial certificate sizes.

McCormick and multilinear envelope formulas are established tools
\citep{mccormick1976Computability,rikun1997MultilinearEnvelope}.
We use them to relate gap-zero faces and graph projections to certificate
counts. Interval overestimation and validated exclusion regions are also
established tools
\citep{schichlMarkotNeumaier2014ExclusionRegions}. Our objective-cutoff
propagation analysis fixes the expression graph and distinguishes its ideal
bound from its round cost. Competitive comparison with an optimal certificate has analogues in
other approximation problems, including analysis of the chord algorithm
\citep{daskalakisDiakonikolasYannakakis2016ChordAlgorithm}. Our branching
results compare tree counts under an exact quadratic-gap oracle, with
explicit information and incumbent assumptions.

Tree decompositions and reuse of conditioned subproblems are established
in discrete AND/OR branch-and-bound and weighted constraint optimization
\citep{marinescuDechter2009AndOrBranchAndBound,deGivrySchiexVerfaillie2006TreeDecomposition}.
Treewidth also controls approximate LP formulations for polynomial
optimization \citep{bienstockMunoz2018LpFormulations}. The focused
decomposition comparison here fixes the continuous factor relaxation and
proves a separation between global-box and affine-separator certificate
sizes. Its existence and checking statements keep the costs of local
optimization explicit.

Integer branch-and-bound lower bounds already use midpoint conflicts and
bounds on the number of integer points that one leaf can contain
\citep{deyDubeyMolinaro2023BranchAndBoundLowerBounds}. Hiding sets give a
related obstruction to small formulations
\citep{kaibelWeltge2015LowerBoundsBranchAndBound}.
The class-number formulation records the full convex-class covering
obstruction for a specified projected relaxation. It also isolates the
additional restriction imposed by variable or split branching. The
statistical applications use these distinctions to compare conflict bounds,
convex-piece certificates, and path certificates.

The Boolean and perspective formulations of sparse regression are prior
work
\citep{pilanciWainwrightElGhaoui2015BooleanRelaxation,xieDeng2018BooleanRelaxation,dongChenLinderoth2015SparseRegression}.
Gaussian support-recovery thresholds for other estimators, including the
Lasso, concern different optimization and certificate events
\citep{wainwright2009SharpLassoThresholds}. Our sparse-regression results
compare root certification, single-variable certification, and selected
local lifts in explicit asymptotic regimes. We also prove a narrow
fixed-dimensional limitation of the Gaussian exactness assertion in
\citet[Theorem~2]{pilanciWainwrightElGhaoui2015BooleanRelaxation} under its
printed per-entry noise and ridge normalization: the true support becomes
optimal, while the exactness probability approaches a constant below one.
The underlying reformulation and deterministic exactness conditions remain
valid.

For binary least squares, maximum-likelihood recovery and box-decoder
recovery have been studied previously
\citep{hansenHassibiDimakisXu2009NearOptimalDetection,huLu2020LimitingPoissonMIMO}.
A theorem about rounding one relaxation supplies no simultaneous conclusion
for all wrong-fixing nodes without an appropriate probability bound.
Gaussian cone face-dimension and chi-square projection mixtures are
classical
\citep{hugSchneider2016RandomConicalTessellations,mccoyTropp2014SteinerFormulas}.
We give a direct sign-orbit proof and use a finite coefficient tail to control
all single-wrong-fixing nodes, including square systems. The resulting C1
threshold, root inactivity threshold, and certificate-size law address
distinct branch-and-bound questions.

\subsection{Scope and organization}\label{intro:scope}

The continuous constants can grow exponentially with dimension, and some
explicit statistical lower-bound constants make the bounds vacuous at the
sizes of the archived computations. \Cref{sec:experiments} presents those
computations as finite-size evidence with their recorded solver settings,
precision, and provenance. They illustrate mechanisms and limitations;
they establish no general practical speedup.

\Cref{sec:certificates} defines the certificate classes and operation
accounting. \Cref{sec:geometry} proves covering, stratified, constraint-gap,
and McCormick results. \Cref{sec:face-exact} gives the root-face integral
characterization and analytic consequences, and \cref{sec:decomposition}
compares spatial and decomposition certificates under the same local
oracle. \Cref{sec:propagation,sec:branching} study propagation and branching.
\Cref{sec:lattice} develops convex integer certificates, followed by sparse
regression and binary least squares in
\cref{sec:regression,sec:binary-least-squares}.
\Cref{sec:experiments,sec:discussion} give the computational evidence,
limitations, and open questions. Long proofs are collected in the
appendices.
